% Generated draft: requires mathematical and source review.
\documentclass[11pt]{article}
\usepackage[margin=25mm]{geometry}
\usepackage{fontspec}
\setmainfont{Latin Modern Roman}
\usepackage{amsmath,amssymb,mathtools,unicode-math}
\setmathfont{Latin Modern Math}
\usepackage{enumitem,needspace,xurl,hyperref}
\setlength{\parindent}{0pt}
\setlength{\parskip}{5pt}
\setlength{\emergencystretch}{3em}
\newcommand{\N}{\mathbb{N}}
\newcommand{\Z}{\mathbb{Z}}
\newcommand{\Q}{\mathbb{Q}}
\newcommand{\R}{\mathbb{R}}
\newcommand{\C}{\mathbb{C}}
\newcommand{\F}{\mathbb{F}}
\newcommand{\A}{\mathbb{A}}
\newcommand{\PP}{\mathbb{P}}
\newcommand{\E}{\mathbb{E}}
\newcommand{\eps}{\varepsilon}
\newcommand{\dd}{\,\mathrm{d}}
\begin{document}
\section*{UnsolvedMath ID: 30003327}
\textbf{Source problem code:} catalogue entry
\subsection{Definitions and mathematical statement}
An ordered exponential field is conventionally a structure \(\mathcal K=(K,0,1,+,-,\cdot,<,\exp_K)\) in which \(K\) is an ordered field and \(\exp_K:(K,+)\to(K_{>0},\cdot)\) is an order-preserving group isomorphism. Work in the language \(L_{\exp}=\{0,1,+,-,\cdot,<,\exp\}\). The structure is o-minimal if every subset of \(K\) definable with parameters in this language is a finite union of points and intervals.

The supplied record additionally requires an axiom or condition called \(\mathrm{EXP}\), but does not state that condition. Accordingly, \(\mathrm{EXP}(\mathcal K)\) below denotes the unspecified condition intended by the source; this draft does not supply a replacement definition.

To describe residue exponentiation, let \(\mathcal O\subseteq K\) be a convex valuation ring with maximal ideal \(\mathfrak m\), and let \(k=\mathcal O/\mathfrak m\), equipped with its induced ordering. An induced residue exponential is given by
\[
\exp_k(a+\mathfrak m)=\exp_K(a)+\mathfrak m\qquad(a\in\mathcal O),
\]
provided this formula is well-defined and makes \(k\) an ordered exponential field. In particular, its use requires that exponentials of elements of \(\mathcal O\) reduce to positive nonzero residue classes and that congruent inputs have congruent exponentials. The record does not specify the valuation ring or explain these compatibility requirements. If the natural valuation is intended, its valuation ring and maximal ideal are
\[
\mathcal O_{\mathrm{nat}}=\{a\in K: |a|\le n\text{ for some positive integer }n\},
\qquad
\mathfrak m_{\mathrm{nat}}=\{a\in K: |a|<1/n\text{ for every positive integer }n\}.
\]
This possible interpretation is recorded explicitly rather than imposed as an additional source hypothesis. Write \(\mathcal K_{\mathrm{res}}\) for the residue exponential structure intended in the supplied record.

The requested assertion is the following formulation schema, pending clarification of \(\mathrm{EXP}\) and the residue construction:
\[
\forall\mathcal K\;\Bigl[
\bigl(\mathcal K\text{ is an o-minimal ordered exponential field}\ \land\ \mathrm{EXP}(\mathcal K)\bigr)
\Longrightarrow
\exists j:\mathcal K_{\mathrm{res}}\hookrightarrow\mathcal K\text{ elementary in }L_{\exp}
\Bigr].
\]
Here elementary means that for every nonnegative integer \(n\), every \(L_{\exp}\)-formula \(\varphi(x_1,\ldots,x_n)\), and every \(\bar a\in K_{\mathrm{res}}^n\),
\[
\mathcal K_{\mathrm{res}}\models\varphi(\bar a)
\quad\Longleftrightarrow\quad
\mathcal K\models\varphi(j(\bar a)).
\]
The question asks for such an embedding into \(\mathcal K\) itself. The source describes this as required by the strongest residue-field conjecture, but does not supply that conjecture's full formulation.
\subsection{Short English statement}
The supplied problem asks whether the residue exponential field of every o-minimal ordered exponential field satisfying \(\mathrm{EXP}\) admits an elementary embedding into the original field, preserving all first-order formulas in the ordered exponential language. The meaning of \(\mathrm{EXP}\) and the intended residue construction require clarification before this becomes a fully specified assertion.
\subsection{Sources}
\begin{itemize}
\item Supplied local UnsolvedMath registry record, permanent ID \(30003327\), title \emph{Elementary Embeddings of Residue Exponential Fields}, source problem code \texttt{catalogue entry}, in the Oberwolfach Reports Open Problems collection. No individual proposer is supplied; the proposed year is recorded as 2016.
\item Supplied report citation: \emph{Mini-Workshop: Surreal Numbers, Surreal Analysis, Hahn Fields and Derivations} (2017), Oberwolfach Reports, pp. 3313--3372, DOI \texttt{10.4171/owr/2016/60}. Original OWR report: \url{https://doi.org/10.4171/owr/2016/60}. The record describes these reports as documenting mathematical workshops and their open-problem discussions.
\item Supplied literature source: Krapp, \emph{Value Groups and Residue Fields}, \url{https://arxiv.org/abs/1803.03153}.
\item Supplied literature source: Fields Institute research discussion, \url{https://www.fields.utoronto.ca/talks/o-minimal-EXP-fields-and-Schanuels-conjecture}.
\item Literature assessment dated exactly \texttt{2026-08-21}: the supplied curation labels the problem open and reports partial progress through residue-field embeddability and related Schanuel-conditional results, while describing the conjectured elementary equivalence and elementary embedding for every o-minimal EXP-field as unresolved. This is supplied dated curation derived from the OWR record, not an independently verified literature assessment; its time-sensitive status and cited claims require independent verification.
\item The registry labels the statement \texttt{corrected\_verified} and reports that the official report passage was checked and rewritten to preserve its mathematical content. These are supplied verification claims, not verification performed for this draft.
\end{itemize}
\end{document}
